%% ma: L12-B14-0009
%% lop: 12
%% chuong: 5
%% bai: 14
%% dang: 12.14.9
%% loai: TN4
%% muc: VD
%% dap_an: "D"
%% nguon: "(NBV)-ĐỀ SỐ 12-MỖI NGÀY 1 ĐỀ THI 2025.pdf (D12, MỖI NGÀY 1 ĐỀ - Đề số 12), Phần I Câu 12"
%% kien_thuc: "Bài 14 (khoảng cách từ điểm đến mặt phẳng), Bài 7 (tọa độ trung điểm); tính chất đối xứng tâm"
%% trang_thai: cho-duyet
%% ly_do: "Dạng toán mới đề xuất, chờ giáo viên duyệt: 12.14.9"
%% ghi_chu: "Đáp án gốc D đúng."
\begin{ex}
Trong không gian $Oxyz$, cho hình bình hành $ABCD$ có tâm $I(2;1;-3)$, đỉnh $C$ nằm trên mặt phẳng $(P)\colon x+2y-3z-1=0$. Khoảng cách từ $A$ đến mặt phẳng $(P)$ bằng
\choice{$\dfrac{24\sqrt{14}}{7}$}{$\dfrac{12\sqrt7}{7}$}{$\dfrac{24\sqrt7}{7}$}{\True $\dfrac{12\sqrt{14}}{7}$}
\loigiai{$I$ là trung điểm của $AC$ và $C\in(P)$, nên khoảng cách từ $A$ đến $(P)$ gấp đôi khoảng cách từ $I$ đến $(P)$ (hai tam giác vuông đồng dạng với tỉ số $2$).

$d(I,(P))=\dfrac{|2+2\cdot1-3\cdot(-3)-1|}{\sqrt{1^2+2^2+(-3)^2}}=\dfrac{12}{\sqrt{14}}$.

Vậy $d(A,(P))=2\cdot\dfrac{12}{\sqrt{14}}=\dfrac{24}{\sqrt{14}}=\dfrac{12\sqrt{14}}{7}$. Chọn D.}
\end{ex}
